% Appendix: development history. Included because several of the errors
% below were invisible to every training and validation metric, and the
% sequence of failures is more instructive than the final configuration.

\section{Development history: what failed, and how it was found}
\label{app:history}

Papers normally report the final configuration. We report the path, because
four of the defects described here produced healthy-looking training curves,
plausible validation losses, and publication-ready figures while being
wrong, and because the method that eventually exposed each one is
transferable. Every stage below corresponds to a dated entry with commit
hashes, shot lists, and metrics in the project's experiment ledger.

\subsection{Stage 1: a latent that could not fail}
\label{app:hist-proxy}

The initial model drove the latent with a relaxation law whose target was
derived from the D$_\alpha$ signal, and predicted D$_\alpha$ through a
parameter-free head equal to $1 - \bact$. With the observation loss weighted
heavily, this pins $\bact \approx 1 - \tilde\yD$ pointwise. Regime
classification from the latent was then classification from D$_\alpha$, and
the model could not be wrong in an interesting way: it was reporting its own
input. Diagnosing this required no experiment, only writing down what the
two equations imply together.

The fix was structural: the latent is driven by actuators alone, and
D$_\alpha$ supervises a \emph{learned} observation head. This is what makes
the regime state falsifiable --- an actuator-driven latent must now predict
the emission collapse rather than copy it, and must decline to predict one
in discharges that do not transition.

\subsection{Stage 2: a redundant parameter reported as physics}
\label{app:hist-cubic}

The latent was first written as
$\tau\dot\zeta = a + c_1\zeta + c_2\zeta^2 - \zeta^3$ and $c_2$ was reported
as an identified coefficient. It is not identifiable: the substitution
$\zeta = \lat + c_2/3$ removes it, and the induced shifts are absorbed by
the drive intercept and the observation offset, which are both free
(Appendix~\ref{app:canonical}). Two parameter sets differing only in $c_2$
describe identical dynamics and identical observations. All subsequent work
uses the depressed cubic, in which $\beta$ is the only topology parameter.

\subsection{Stage 3: loss weights chosen to make a term ``matter''}
\label{app:hist-weights}

Early configurations set the regime and observation weights by inspecting
the magnitude of the profile term ($\mathcal{O}(700)$ in the then-current
units) and raising the others until they were ``visible''. This is
scale-chasing, not modelling: it made the weights uninterpretable and
coupled them to arbitrary unit choices. The objective was rewritten so each
term is an average of order-one quantities with fixed robust scales
(\S\ref{sec:ident-loss}), after which the weights became ratios that can be
stated and varied meaningfully.

\subsection{Stage 4: a finite-volume scheme that did not converge}
\label{app:hist-fvm}

The discretisation used a face-averaged $V'$ for cell volumes and gave the
axis cell full width, both inconsistent with the face positions used in the
flux stencil. The scheme remained stable, trained normally, and produced
sensible fits. Its error under grid refinement was $5.4, 8.1, 9.7,
10.0 \times 10^{-3}$ for $N = 17, 33, 65, 129$ --- an apparent convergence
order of $-0.31$, i.e.\ refining the grid made it worse. No training or
validation metric was sensitive to this. It was found by writing a
manufactured-solution test with an analytic decay mode and refining.

A secondary lesson: the first version of that test used a unit-amplitude
solution and appeared to show first-order convergence stalling at
$10^{-3}$. The solver's smooth positivity safeguard is exact only to
$\sim$$\SI{0.01}{eV}$, so a unit-amplitude test measures the safeguard, not
the discretisation. Re-running at $\SI{1000}{eV}$ --- the operating
magnitude --- gave clean first-order convergence. Verification tests must be
posed at the scale of the quantity they verify.

\subsection{Stage 5: an inverted radial coordinate}
\label{app:hist-rho}

The most consequential error. MAST Level-2 stores poloidal flux per radian,
maximal on the magnetic axis; the normalisation assumed it minimal there and
took grid extrema as the anchors, with a guard clause that discarded the
separatrix value whenever it lay below the grid maximum --- which it always
does. The resulting coordinate had $\rho = 1$ \emph{at} the magnetic axis
and $\rho \approx 0.75$--$0.82$ at both ends of the Thomson chord,
non-monotonic in major radius. Every quality gate and mask named ``edge''
therefore selected the near-axis region, and the interpolation interleaved
inboard and outboard channels that mapped to similar $\rho$.

Nothing in the training or validation pipeline was sensitive to this either.
The model happily fitted the region it was given. It surfaced only when the
equilibrium data path was audited line by line against the actual file
contents in order to document it --- that is, the error was found by the act
of writing the methods section, not by any diagnostic.

Correcting it (Appendix~\ref{app:eq-rho}) reduced held-out profile error by
a factor $2.6$, moved the best checkpoint from step $150$--$250$ to
$\approx 1850$, and reversed the study's headline conclusion
(\S\ref{sec:res-modelcomp}). We consider the checkpoint shift the most
useful transferable diagnostic here: a model that plateaus almost
immediately may not be converged --- it may have nothing real to learn.

\subsection{Stage 6: geometry that was never the reconstruction's}
\label{app:hist-vprime}

The metric $V'(\rho)$ was documented as coming from the equilibrium. In fact
the Level-2 group contains no flux-surface volume profile, so the code fell
through a sentinel path to the analytic cylindrical $V' = 2\rho$ for every
discharge. The claim was false, and the sentinel was value-based
($V' \equiv 1$), so a genuine near-unity metric would have been silently
overwritten too. We now integrate $V(\rho)$ from the flux map directly and
validate it against the equilibrium's own total-volume scalar, agreeing to
$0.1$--$0.3\%$ (Appendix~\ref{app:eq-vprime}). The general lesson is that a
fallback which produces plausible numbers is more dangerous than one that
crashes.

\subsection{Stage 7: probabilities that were never calibrated}
\label{app:hist-calib}

The model reported $\pH$ as a probability and was evaluated with
threshold-dependent scores. Its AUC was $0.84$--$0.94$ while $F_1$ was
exactly zero: discrimination was good, but the raw probabilities almost
never crossed $0.5$, so no positive prediction was ever made. An earlier
draft had reported a median transition-timing error computed only over the
minority of discharges where a crossing happened to occur --- a selected
subset, and therefore a meaningless statistic. Calibration is now measured
(reliability curve, expected calibration error, Brier score), performed by
Platt scaling fitted on validation data only, and reported alongside the
uncalibrated values (\S\ref{sec:ident-select}).

\subsection{Stage 8: diagnosing a failure instead of tuning around it}
\label{app:hist-covariates}

With the above fixed, L$\to$H event recall was $0.06$--$0.11$: poor. The
tempting responses --- more capacity, longer training, different weights ---
would all have been wrong. Inspecting the latent trajectories showed the
model was behaving correctly: in the discharges it missed, the identified
drive never reached the fold, and the latent therefore stayed on the lower
branch, as it should. Those discharges were concentrated in the X-point-height
threshold campaigns, whose transition threshold depends on plasma shape ---
a quantity the drive could not see.

Extending the drive with a loss-power proxy and shape covariates raised the
fraction of discharges in which the latent switches branches from $\sim$$30\%$
to $100\%$ and roughly quadrupled event recall. The methodological point is
that a performance deficit was traced to a specific missing physical
covariate by examining the model's internal state, rather than treated as an
optimisation problem.

\subsection{Stage 9: protocol errors}
\label{app:hist-protocol}

Two process failures are recorded for completeness. First, an early
comparison ran with an early-stopping patience of eight evaluations while
the typical interval between validation improvements was of the same order,
so some seeds stopped while still improving; the patience was doubled and
all affected runs were discarded and rerun rather than mixed. Second, the
training script persisted the raw rather than the override-applied
configuration, so a constrained-topology run was first evaluated under the
wrong parameterisation; this was caught by the identified $\beta$ having the
wrong sign for its arm, and both the bug and the re-evaluation are recorded.

\subsection{What this history implies for the results}
\label{app:hist-implications}

Three of the defects above (Stages 4, 5, 6) were invisible to every loss and
metric we computed, and two of them changed conclusions. The safeguards that
actually caught them were: a manufactured-solution convergence study, a
line-by-line audit of the data path against file contents, and an
independent physical validation of a derived quantity against a value the
files already contained. None of these is a machine-learning diagnostic. We
would suggest that any study identifying physical parameters from
experimental data through a differentiable solver should budget for all
three before interpreting a fitted coefficient.
